% Môn: Vật lý
% Lớp: 12
% Chương: Chương II. Khí lí tưởng
% Bài: L12C2 Bài 10. Định luật Charles · Ứng dụng định luật Boyle trong bơm, xi lanh và đời sống
% ===== Câu 54 =====
% ID: L12C2B10-218-DS
% Mức: VD
\dangbt{Ứng dụng định luật Boyle trong bơm, xi lanh và đời sống}
\begin{ex}
% ID: L12C2B10-218-DS
% Mức: VD
Trong dạng “Nhận biết quá trình đẳng nhiệt và định luật Boyle”, Giải thích hoạt động của bơm tiêm khi bịt đầu và đẩy pít-tông chậm.
\choiceTF

\loigiai{
Thể tích khí giảm trong điều kiện gần đẳng nhiệt nên áp suất tăng theo định luật Boyle, tạo lực cản lên pít-tông.
}
\end{ex}

% ===== Câu 78 =====
% ID: L12C2B10-219-DS
% Mức: VD
\begin{ex}
% ID: L12C2B10-219-DS
% Mức: VD
Trong dạng “Tính áp suất, thể tích trong quá trình đẳng nhiệt”, Giải thích hoạt động của bơm tiêm khi bịt đầu và đẩy pít-tông chậm.
\choiceTF

\loigiai{
Thể tích khí giảm trong điều kiện gần đẳng nhiệt nên áp suất tăng theo định luật Boyle, tạo lực cản lên pít-tông.
}
\end{ex}

% ===== Câu 95 =====
% ID: L12C2B10-220-TLN
% Mức: VD
\begin{ex}
% ID: L12C2B10-220-TLN
% Mức: VD
\shortans{5}
\loigiai{
Thể tích ban đầu của khí là $V_1 = S \cdot L_1$, trong đó $S$ là diện tích đáy xi lanh và $L_1$ là chiều dài ban đầu của cột khí.
Diện tích đáy xi lanh là $S = \pi \cdot (D/2)^2 = \pi \cdot (10/2)^2 = \pi \cdot 5^2 = 25\pi$ cm$^2$.
Thể tích ban đầu $V_1 = 25\pi \cdot 20 = 500\pi$ cm$^3$.
Vì quá trình nén là đẳng nhiệt, ta áp dụng định luật Boyle-Mariotte:
$P_1 \cdot V_1 = P_2 \cdot V_2$
Trong đó:
$P_1 = 1,2$ atm
$V_1 = 500\pi$ cm$^3P_2 = 4,8$ atm
$V_2 = S \cdot L_2 = 25\pi \cdot L_2$
Thay các giá trị vào phương trình:
$1,2 \cdot 500\pi = 4,8 \cdot 25\pi \cdot L_2$
Chia cả hai vế cho $25\pi$:
$1,2 \cdot \frac{500\pi}{25\pi} = 4,8 \cdot L_21,2 \cdot 20 = 4,8 \cdot L_224 = 4,8 \cdot L_2$
Tính $L_2$:
$L_2 = \frac{24}{4,8} = 5$ cm.
Vậy, chiều dài cuối cùng của cột khí là 5 cm.
}
\end{ex}

% ===== Câu 102 =====
% ID: L12C2B10-221-DS
% Mức: VD
\begin{ex}
% ID: L12C2B10-221-DS
% Mức: VD
Trong dạng “Đồ thị đẳng nhiệt trên các hệ tọa độ”, Giải thích hoạt động của bơm tiêm khi bịt đầu và đẩy pít-tông chậm.
\choiceTF

\loigiai{
Thể tích khí giảm trong điều kiện gần đẳng nhiệt nên áp suất tăng theo định luật Boyle, tạo lực cản lên pít-tông.
}
\end{ex}

% ===== Câu 113 =====
% ID: L12C2B10-222-DS
% Mức: TH
\begin{ex}
% ID: L12C2B10-222-DS
% Mức: TH
Xét các phát biểu sau trong dạng “Ứng dụng định luật Boyle trong bơm, xi lanh và đời sống”.
\choiceTF
{\True Đại lượng nào không đổi trong quá trình đẳng nhiệt của một lượng khí xác định? — phát biểu đúng là: Nhiệt độ.}
{Đại lượng nào không đổi trong quá trình đẳng nhiệt của một lượng khí xác định? — phát biểu Áp suất là đúng.}
{\True Hệ thức nào là định luật Boyle? — phát biểu đúng là: $p_1V_1=p_2V_2$.}
{Hệ thức nào là định luật Boyle? — phát biểu $V_1/T_1=V_2/T_2$ là đúng.}
\loigiai{
\begin{itemchoice}
\item (Đúng) Đại lượng nào không đổi trong quá trình đẳng nhiệt của một lượng khí xác định? — phát biểu đúng là: Nhiệt độ. (Vì): Quá trình đẳng nhiệt có nhiệt độ không đổi. b) Sai: trái với hệ thức hoặc điều kiện đã nêu. c) Đúng: Ở nhiệt độ không đổi, $pV$ là hằng số. d) Sai: trái với hệ thức hoặc điều kiện đã nêu
\item (Sai) Đại lượng nào không đổi trong quá trình đẳng nhiệt của một lượng khí xác định? — phát biểu Áp suất là đúng.
\item (Đúng) Hệ thức nào là định luật Boyle? — phát biểu đúng là: $p_1V_1=p_2V_2$.
\item (Sai) Hệ thức nào là định luật Boyle? — phát biểu $V_1/T_1=V_2/T_2$ là đúng.
\end{itemchoice}
}
\end{ex}

% ===== Câu 114 =====
% ID: L12C2B10-223-DS
% Mức: TH
\begin{ex}
% ID: L12C2B10-223-DS
% Mức: TH
Trong dạng “Ứng dụng định luật Boyle trong bơm, xi lanh và đời sống”, Trình bày cơ sở lí thuyết của dạng “Ứng dụng định luật Boyle trong bơm, xi lanh và đời sống” và nêu điều kiện áp dụng.
\choiceTF

\loigiai{
Nêu đúng đại lượng không đổi, hệ thức đặc trưng và yêu cầu dùng nhiệt độ kelvin, áp suất tuyệt đối khi có liên quan.
}
\end{ex}

% ===== Câu 116 =====
% ID: L12C2B10-224-TN
% Mức: VD
\begin{ex}
% ID: L12C2B10-224-TN
% Mức: VD
Trong dạng “Nhận biết quá trình đẳng nhiệt và định luật Boyle”, Bơm xe đạp nóng lên rõ rệt thì quá trình nén
\choice
{\True không hoàn toàn đẳng nhiệt}
{luôn đẳng nhiệt chính xác}
{đẳng áp}
{đẳng tích}
\loigiai{
Đáp án A. Nhiệt độ thay đổi nên không thể áp dụng Boyle chính xác như quá trình đẳng nhiệt.
}
\end{ex}

% ===== Câu 117 =====
% ID: L12C2B10-225-DS
% Mức: TH
\begin{ex}
% ID: L12C2B10-225-DS
% Mức: TH
Xét các phát biểu sau trong dạng “Ứng dụng định luật Boyle trong bơm, xi lanh và đời sống”.
\choiceTF
{\True Khí có áp suất $125\,\text{kPa}$, thể tích 20 lít được nén đẳng nhiệt còn 4 lít. Áp suất sau là — phát biểu đúng là: $625\,\text{kPa}$.}
{Khí có áp suất $125\,\text{kPa}$, thể tích 20 lít được nén đẳng nhiệt còn 4 lít. Áp suất sau là — phát biểu $25\,\text{kPa}$ là đúng.}
{\True Khí có áp suất $130\,\text{kPa}$, thể tích 6 lít giãn đẳng nhiệt đến 12 lít. Áp suất sau là — phát biểu đúng là: $65\,\text{kPa}$.}
{Khí có áp suất $130\,\text{kPa}$, thể tích 6 lít giãn đẳng nhiệt đến 12 lít. Áp suất sau là — phát biểu $130\,\text{kPa}$ là đúng.}
\loigiai{
\begin{itemchoice}
\item (Đúng) Khí có áp suất $125\,\text{kPa}$, thể tích 20 lít được nén đẳng nhiệt còn 4 lít. Áp suất sau là — phát biểu đúng là: $625\,\text{kPa}$. (Vì): $p_2=p_1V_1/V_2=125\cdot20/4=625$ kPa. b) Sai: trái với hệ thức hoặc điều kiện đã nêu. c) Đúng: $p_2=130\cdot6/12=65$ kPa. d) Sai: trái với hệ thức hoặc điều kiện đã nêu
\item (Sai) Khí có áp suất $125\,\text{kPa}$, thể tích 20 lít được nén đẳng nhiệt còn 4 lít. Áp suất sau là — phát biểu $25\,\text{kPa}$ là đúng.
\item (Đúng) Khí có áp suất $130\,\text{kPa}$, thể tích 6 lít giãn đẳng nhiệt đến 12 lít. Áp suất sau là — phát biểu đúng là: $65\,\text{kPa}$.
\item (Sai) Khí có áp suất $130\,\text{kPa}$, thể tích 6 lít giãn đẳng nhiệt đến 12 lít. Áp suất sau là — phát biểu $130\,\text{kPa}$ là đúng.
\end{itemchoice}
}
\end{ex}

% ===== Câu 118 =====
% ID: L12C2B10-226-DS
% Mức: TH
\begin{ex}
% ID: L12C2B10-226-DS
% Mức: TH
Trong dạng “Ứng dụng định luật Boyle trong bơm, xi lanh và đời sống”, Mô tả dạng đồ thị đặc trưng và cách nhận biết quá trình trong dạng “Ứng dụng định luật Boyle trong bơm, xi lanh và đời sống”.
\choiceTF

\loigiai{
Xác định các trục, đại lượng không đổi và suy ra hình dạng đồ thị từ hệ thức vật lí tương ứng.
}
\end{ex}

% ===== Câu 120 =====
% ID: L12C2B10-227-TN
% Mức: VDC
\begin{ex}
% ID: L12C2B10-227-TN
% Mức: VDC
Trong dạng “Nhận biết quá trình đẳng nhiệt và định luật Boyle”, Để quá trình gần đẳng nhiệt nên nén khí
\choice
{\True chậm để trao đổi nhiệt}
{thật nhanh}
{trong bình cách nhiệt}
{với thể tích không đổi}
\loigiai{
Đáp án A. Nén chậm giúp khí trao đổi nhiệt và giữ nhiệt độ gần không đổi.
}
\end{ex}

% ===== Câu 125 =====
% ID: L12C2B10-228-DS
% Mức: VD
\begin{ex}
% ID: L12C2B10-228-DS
% Mức: VD
Xét các phát biểu sau trong dạng “Ứng dụng định luật Boyle trong bơm, xi lanh và đời sống”.
\choiceTF
{\True Khi dùng định luật Boyle phải sử dụng — phát biểu đúng là: áp suất tuyệt đối.}
{Khi dùng định luật Boyle phải sử dụng — phát biểu áp suất dư là đúng.}
{\True Muốn nén khí gần đẳng nhiệt nên thực hiện — phát biểu đúng là: chậm để khí trao đổi nhiệt.}
{Muốn nén khí gần đẳng nhiệt nên thực hiện — phát biểu ở thể tích không đổi là đúng.}
\loigiai{
\begin{itemchoice}
\item (Đúng) Khi dùng định luật Boyle phải sử dụng — phát biểu đúng là: áp suất tuyệt đối. (Vì): Phương trình trạng thái dùng áp suất tuyệt đối của khí. b) Sai: trái với hệ thức hoặc điều kiện đã nêu. c) Đúng: Nén chậm giúp nhiệt độ khí gần như không đổi. d) Sai: trái với hệ thức hoặc điều kiện đã nêu
\item (Sai) Khi dùng định luật Boyle phải sử dụng — phát biểu áp suất dư là đúng.
\item (Đúng) Muốn nén khí gần đẳng nhiệt nên thực hiện — phát biểu đúng là: chậm để khí trao đổi nhiệt.
\item (Sai) Muốn nén khí gần đẳng nhiệt nên thực hiện — phát biểu ở thể tích không đổi là đúng.
\end{itemchoice}
}
\end{ex}

% ===== Câu 129 =====
% ID: L12C2B10-229-DS
% Mức: VDC
\begin{ex}
% ID: L12C2B10-229-DS
% Mức: VDC
Xét các phát biểu sau trong dạng “Ứng dụng định luật Boyle trong bơm, xi lanh và đời sống”.
\choiceTF
{\True Với lượng khí cố định đẳng nhiệt, khối lượng riêng tỉ lệ — phát biểu đúng là: thuận với áp suất.}
{Với lượng khí cố định đẳng nhiệt, khối lượng riêng tỉ lệ — phát biểu nghịch với áp suất là đúng.}
{\True Trong ống có cột chất lỏng, áp suất khí bị giam được xác định bằng — phát biểu đúng là: cân bằng áp suất tại cùng độ cao.}
{Trong ống có cột chất lỏng, áp suất khí bị giam được xác định bằng — phát biểu chỉ $\rho gh$ là đúng.}
\loigiai{
\begin{itemchoice}
\item (Đúng) Với lượng khí cố định đẳng nhiệt, khối lượng riêng tỉ lệ — phát biểu đúng là: thuận với áp suất. (Vì): Vì $\rho=m/V$ và $pV$ không đổi nên $\rho\sim p$. b) Sai: trái với hệ thức hoặc điều kiện đã nêu. c) Đúng: Dùng $p_{khí}=p_0\pm\rho gh$ tùy cách bố trí. d) Sai: trái với hệ thức hoặc điều kiện đã nêu
\item (Sai) Với lượng khí cố định đẳng nhiệt, khối lượng riêng tỉ lệ — phát biểu nghịch với áp suất là đúng.
\item (Đúng) Trong ống có cột chất lỏng, áp suất khí bị giam được xác định bằng — phát biểu đúng là: cân bằng áp suất tại cùng độ cao.
\item (Sai) Trong ống có cột chất lỏng, áp suất khí bị giam được xác định bằng — phát biểu chỉ $\rho gh$ là đúng.
\end{itemchoice}
}
\end{ex}

% ===== Câu 133 =====
% ID: L12C2B10-230-TL
% Mức: VDC
\begin{bt}
% ID: L12C2B10-230-TL
% Mức: VDC
Trong dạng “Ứng dụng định luật Boyle trong bơm, xi lanh và đời sống”, undefined
\loigiai{
$V_2=140\cdot18/(420)=6$ lít.
}
\end{bt}

% ===== Câu 142 =====
% ID: L12C2B10-231-TN
% Mức: VD
\begin{ex}
% ID: L12C2B10-231-TN
% Mức: VD
Trong dạng “Tính áp suất, thể tích trong quá trình đẳng nhiệt”, Bơm xe đạp nóng lên rõ rệt thì quá trình nén
\choice
{\True không hoàn toàn đẳng nhiệt}
{luôn đẳng nhiệt chính xác}
{đẳng áp}
{đẳng tích}
\loigiai{
Đáp án A. Nhiệt độ thay đổi nên không thể áp dụng Boyle chính xác như quá trình đẳng nhiệt.
}
\end{ex}

% ===== Câu 143 =====
% ID: L12C2B10-232-TN
% Mức: VDC
\begin{ex}
% ID: L12C2B10-232-TN
% Mức: VDC
Trong dạng “Tính áp suất, thể tích trong quá trình đẳng nhiệt”, Để quá trình gần đẳng nhiệt nên nén khí
\choice
{\True chậm để trao đổi nhiệt}
{thật nhanh}
{trong bình cách nhiệt}
{với thể tích không đổi}
\loigiai{
Đáp án A. Nén chậm giúp khí trao đổi nhiệt và giữ nhiệt độ gần không đổi.
}
\end{ex}

% ===== Câu 152 =====
% ID: L12C2B10-233-TN
% Mức: VD
\begin{ex}
% ID: L12C2B10-233-TN
% Mức: VD
Trong dạng “Đồ thị đẳng nhiệt trên các hệ tọa độ”, Bơm xe đạp nóng lên rõ rệt thì quá trình nén
\choice
{\True không hoàn toàn đẳng nhiệt}
{luôn đẳng nhiệt chính xác}
{đẳng áp}
{đẳng tích}
\loigiai{
Đáp án A. Nhiệt độ thay đổi nên không thể áp dụng Boyle chính xác như quá trình đẳng nhiệt.
}
\end{ex}

% ===== Câu 153 =====
% ID: L12C2B10-234-TN
% Mức: VDC
\begin{ex}
% ID: L12C2B10-234-TN
% Mức: VDC
Trong dạng “Đồ thị đẳng nhiệt trên các hệ tọa độ”, Để quá trình gần đẳng nhiệt nên nén khí
\choice
{\True chậm để trao đổi nhiệt}
{thật nhanh}
{trong bình cách nhiệt}
{với thể tích không đổi}
\loigiai{
Đáp án A. Nén chậm giúp khí trao đổi nhiệt và giữ nhiệt độ gần không đổi.
}
\end{ex}
